% GENERATED FILE. Edit canonical lessons, then run npm run generate:books.
% canonical-source-hash: fnv1a64:78f1570dd585bff1
% canonical-lessons: ZH-C398-taidu, ZH-C398-qubie, ZH-C398-diqu, ZH-C398-xiancheng, ZH-C398-qujueyu

\chapter{It Depends}
\label{ch:zh-it-depends}
\hlchaptermodality{398}

\begin{chapteropening}
\textbf{By the end of this chapter:} \emph{I can talk about attitude, the difference between two things, a region, ready-made food and what a plan depends on.}

Complete the last lesson of chapter 398: I can talk about attitude, the difference between two things, a region, ready-made food and what a plan depends on.
\end{chapteropening}

\section[tàidu]{tàidu — attitude}

\label{lesson:ZH-C398-taidu}



\begin{quote}
Before the new one: say the Mandarin for in an awkward spot, then the Mandarin for an excuse (a poor reason).
\end{quote}

\subsection*{You'll want to know: tàidu}
\emph{tàidu} — \textquotedblleft{}attitude\textquotedblright{}.

\begin{quote}
\emph{néng bu néng chénggōng yào kàn nǐ de tàidu} — Whether you succeed depends on your attitude.
\end{quote}

Its characters come later; for now it is a word for the ear.

\begin{cousinweb}
\emph{tàidu} is built from \emph{tài} 'a bearing' and \emph{dù} 'a measure', both inherited from Middle Chinese, \emph{thojH} and \emph{duH} (Guangyun, in Baxter's transcription), which glosses the first 'manner' and the second 'a rule'. In the word the second loses its tone.
\end{cousinweb}

\begin{grammarlens}[title={What you've built}]
One more word for saying what would happen.
\end{grammarlens}

\subsection*{Your turn}
\begin{itemize}
\raggedright
  \item \emph{Say it:} \emph{tàidu}
  \item \emph{Say it:} \emph{tàidu}, once more
  \item \emph{Say it:} \emph{néng bu néng chénggōng yào kàn nǐ de tàidu} — Whether you succeed depends on your attitude.
  \item \emph{Recall:} say the Mandarin for in an awkward spot, then the Mandarin for an excuse (a poor reason), then say \emph{tàidu} again
\end{itemize}

\begin{tcolorbox}[breakable,colback=teal!4,colframe=teal!35!black,title={Before you move on}]
What does tàidu mean? (\textquotedblleft{}Attitude\textquotedblright{}.) Say it once more.
\end{tcolorbox}
\section[qūbié]{qūbié — difference}

\label{lesson:ZH-C398-qubie}



\begin{quote}
Before the new one: say the Mandarin for an excuse (a poor reason), then the Mandarin for your attitude.
\end{quote}

\subsection*{You'll want to know: qūbié}
\emph{qūbié} — \textquotedblleft{}difference\textquotedblright{}.

\begin{quote}
\emph{zhè liǎng ge yǒu shénme qūbié} — What's the difference between these two?
\end{quote}

Its characters come later; for now it is a word for the ear.

\begin{cousinweb}
\emph{qūbié} is built from \zh{区} 'an area' and \emph{bié} 'to part', both inherited from Middle Chinese, \emph{khju} and \emph{bjet} (Guangyun, in Baxter's transcription), which glosses the second 'different' and 'to leave'.
\end{cousinweb}

\begin{grammarlens}[title={What you've built}]
One more word for saying what would happen.
\end{grammarlens}

\subsection*{Your turn}
\begin{itemize}
\raggedright
  \item \emph{Say it:} \emph{qūbié}
  \item \emph{Say it:} \emph{qūbié}, once more
  \item \emph{Say it:} \emph{zhè liǎng ge yǒu shénme qūbié} — What's the difference between these two?
  \item \emph{Recall:} say the Mandarin for an excuse (a poor reason), then the Mandarin for your attitude, then say \emph{qūbié} again
\end{itemize}

\begin{tcolorbox}[breakable,colback=teal!4,colframe=teal!35!black,title={Before you move on}]
What does qūbié mean? (\textquotedblleft{}Difference\textquotedblright{}.) Say it once more.
\end{tcolorbox}
\section[\texorpdfstring{\zh{地区} (dìqū)}{地区 (dìqū)}]{\zh{地区} (dìqū) — region; area (of a country)}

\label{lesson:ZH-C398-diqu}



\begin{quote}
Before the new one: say the Mandarin for your attitude, then the Mandarin for the difference.
\end{quote}

\subsection*{You'll want to know: \zh{地区}}
\textbf{\zh{地区}} — \emph{dìqū} — \textquotedblleft{}region; area (of a country)\textquotedblright{}.

\begin{quote}
\textbf{\zh{这个地区有一百万人口}} \emph{zhège dìqū yǒu yìbǎi wàn rénkǒu} — This region has a million people.
\end{quote}

\begin{cousinweb}
\textbf{\zh{地区}} (\emph{dìqū}) is built from \zh{地} 'the ground' and \zh{区} 'an area', both inherited from Middle Chinese, \emph{dijH} and \emph{khju} (Guangyun, in Baxter's transcription), which glosses the first 'the land'.
\end{cousinweb}

\begin{grammarlens}[title={What you've built}]
One more word for saying what would happen.
\end{grammarlens}

\subsection*{Your turn}
\begin{itemize}
\raggedright
  \item \emph{Say it:} \emph{dìqū}
  \item \emph{Say it:} \emph{dìqū}, once more
  \item \emph{Say it:} \emph{zhège dìqū yǒu yìbǎi wàn rénkǒu} — This region has a million people.
  \item \emph{Recall:} say the Mandarin for your attitude, then the Mandarin for the difference, then say \emph{dìqū} again
\end{itemize}

\begin{tcolorbox}[breakable,colback=teal!4,colframe=teal!35!black,title={Before you move on}]
What does \zh{地区} mean? (\textquotedblleft{}Region; area (of a country)\textquotedblright{}.) Say it once more.
\end{tcolorbox}
\section[\texorpdfstring{\zh{现成} (xiànchéng)}{现成 (xiànchéng)}]{\zh{现成} (xiànchéng) — ready-made}

\label{lesson:ZH-C398-xiancheng}



\begin{quote}
Before the new one: say the Mandarin for the difference, then the Mandarin for a region of a country.
\end{quote}

\subsection*{You'll want to know: \zh{现成}}
\textbf{\zh{现成}} — \emph{xiànchéng} — \textquotedblleft{}ready-made\textquotedblright{}.

\begin{quote}
\textbf{\zh{要是没时间做饭就买现成的}} \emph{yàoshi méi shíjiān zuòfàn jiù mǎi xiànchéng de} — If there's no time to cook, buy something ready-made.
\end{quote}

\begin{cousinweb}
\textbf{\zh{现成}} (\emph{xiànchéng}) is built from \zh{现} 'to appear' and \zh{成} 'complete', both inherited from Middle Chinese, \emph{henH} and \emph{dzyeng} (Guangyun, in Baxter's transcription), which glosses the first 'to show' and the second 'finished': already there and done.
\end{cousinweb}

\begin{grammarlens}[title={What you've built}]
One more word for saying what would happen.
\end{grammarlens}

\subsection*{Your turn}
\begin{itemize}
\raggedright
  \item \emph{Say it:} \emph{xiànchéng}
  \item \emph{Say it:} \emph{xiànchéng}, once more
  \item \emph{Say it:} \emph{yàoshi méi shíjiān zuòfàn jiù mǎi xiànchéng de} — If there's no time to cook, buy something ready-made.
  \item \emph{Recall:} say the Mandarin for the difference, then the Mandarin for a region of a country, then say \emph{xiànchéng} again
\end{itemize}

\begin{tcolorbox}[breakable,colback=teal!4,colframe=teal!35!black,title={Before you move on}]
What does \zh{现成} mean? (\textquotedblleft{}Ready-made\textquotedblright{}.) Say it once more.
\end{tcolorbox}
\section[qǔjuéyú]{qǔjuéyú — to depend on (a factor); to hinge on}

\label{lesson:ZH-C398-qujueyu}



\begin{quote}
Before the new one: say the Mandarin for a region of a country, then the Mandarin for ready-made.
\end{quote}

\subsection*{You'll want to know: qǔjuéyú}
\emph{qǔjuéyú} — \textquotedblleft{}to depend on (a factor); to hinge on\textquotedblright{}.

\begin{quote}
\emph{wǒmen qù bu qù qǔjuéyú tiānqì} — Whether we go hinges on the weather.
\end{quote}

Its characters come later; for now it is a word for the ear.

\begin{cousinweb}
\emph{qǔjuéyú} is built from \zh{取} 'to take', \emph{jué} 'to decide' and \emph{yú} 'by', all inherited from Middle Chinese, \emph{tshjuX}, \emph{kwet} and \emph{'jo} (Guangyun, in Baxter's transcription), which glosses the second 'to break through' and 'to cut' and lists the third as a grammar word: to take its decision from something.
\end{cousinweb}

\begin{grammarlens}[title={What you've built}]
One more word for saying what would happen.
\end{grammarlens}

\subsection*{Your turn}
\begin{itemize}
\raggedright
  \item \emph{Say it:} \emph{qǔjuéyú}
  \item \emph{Say it:} \emph{qǔjuéyú}, once more
  \item \emph{Say it:} \emph{wǒmen qù bu qù qǔjuéyú tiānqì} — Whether we go hinges on the weather.
  \item \emph{Recall:} say the Mandarin for a region of a country, then the Mandarin for ready-made, then say \emph{qǔjuéyú} again
\end{itemize}

\begin{tcolorbox}[breakable,colback=teal!4,colframe=teal!35!black,title={Before you move on}]
What does qǔjuéyú mean? (\textquotedblleft{}To depend on (a factor); to hinge on\textquotedblright{}.) Say it once more.
\end{tcolorbox}
